\section{Introduction}

Digital misinformation often reuses genuine media outside its original context. A photograph from one event may be paired with a claim about another place, date, or incident. This form of out-of-context misinformation is difficult to detect from pixels alone because the image may contain no manipulation artifacts. Pairwise image and text compatibility can reveal a mismatch, but it cannot establish what the image originally depicted or whether the accompanying claim is true. NewsCLIPpings and VERITE illustrate the need to evaluate false context separately from synthetic or visibly edited media \cite{luo2021newsclippings,papadopoulos2024verite}. A fact-checking system therefore needs external evidence that links the claim and image to real events.

Retrieval supplies that context, but the validity of a decision depends on when the evidence appeared and where it originated. For a claim assessed at cutoff time $\tau$, a later correction, official report, or fact-check page contains information that was unavailable at the time of the decision. Allowing such documents into evaluation leaks future information and can inflate measured performance. Search results can also give a false impression of corroboration. Syndicated articles, reposts, and quotation chains may produce many documents that trace back to one report. Document-level aggregation can therefore overweight a single underlying source.

Evidence availability is different from model confidence. The task uses three verdict labels. \emph{Supported} means that time-valid evidence corroborates the claim, \emph{refuted} means that it contradicts the claim, and \emph{insufficient evidence} means that the available material cannot settle the claim. Selective abstention is a separate model decision: the system may defer any prediction when its calibrated confidence is too low. This separation permits independent measurement of evidential sufficiency and decision reliability.

Recent systems address parts of this setting. MOCHEG connects multimodal retrieval with verification \cite{yao2023mocheg}; RED-DOT reranks web evidence by relevance \cite{papadopoulos2025reddot}; and COVE recovers an image's original context before assessing its caption \cite{tonglet2025cove}. KAMP incorporates background knowledge graphs \cite{zhang2025kamp}, while RAGAR and MultiReflect use retrieval-augmented multimodal reasoning \cite{khaliq2024ragar,kangur2025multireflect}. The present study places four concerns in the same decision pipeline: temporal eligibility, source dependence, evidence sufficiency, and calibrated deferral.

This study asks whether time-valid evidence from independent source clusters improves future-event misinformation detection and reduces overconfident decisions relative to ordinary multimodal RAG. Each case is an English claim about a United States event, accompanied by an image and a decision cutoff. \systemname{} filters evidence by publication time, groups dependent sources, represents provenance in a heterogeneous graph, and estimates evidence sufficiency before veracity. Its output includes a verdict, calibrated class probabilities, the selected evidence identifiers, source-cluster identifiers, and the cutoff used for the decision.

The study makes three contributions:
\begin{itemize}
    \item a construction and audit protocol for \datasetname{}, a temporally partitioned benchmark derived from professional fact-check metadata and open web archives, with events, near-duplicate claims, and reused images isolated by split;
    \item a temporal-provenance framework that combines multimodal retrieval, cutoff filtering, source-dependency clustering, graph-based evidence aggregation, hierarchical sufficiency and veracity prediction, and calibrated abstention; and
    \item a controlled evaluation protocol covering future-event generalization, temporal leakage, duplicated-source sensitivity, retrieval quality, calibration, selective risk, and event-level uncertainty.
\end{itemize}

